\chapter{Đào sâu: Cấu hình --- Config Server, Secret và refresh ở quy mô lớn}
\label{ch:dd-config}

Chương~2 giới thiệu config-server như ``thêm một nguồn thuộc tính trong
chồng'': một service HTTP nhỏ trao cho mỗi ứng dụng file YAML của nó lúc
khởi động. Mô tả đó đầy đủ cho một chiếc laptop và thiếu sót nguy hiểm
cho bất kỳ thứ gì khác. Chương này tháo rời đường đi cấu hình của dự án
--- server thực sự phục vụ gì, ai đọc được, chuyện gì xảy ra khi nó thay
đổi, và chuyện gì xảy ra khi nó không có mặt --- rồi xây lại theo hai
cách khác nhau.

Ba câu hỏi sắp xếp chương này:

\begin{enumerate}
  \item \textbf{Cấu hình tập trung giải quyết vấn đề gì}, và trong ba
  backend của nó (native, git, nền tảng) cái nào giải quyết tốt nhất cho
  hệ thống nào?
  \item \textbf{Bán kính thiệt hại của một file cấu hình là gì} --- ai
  đọc nó, quy tắc ưu tiên nào áp dụng, và một giá trị mặc định trong
  file được git theo dõi thực sự có nghĩa gì?
  \item \textbf{Cái gì gãy} khi server chết, khi một giá trị phải đổi
  lúc chạy, và khi một secret phải xoay?
\end{enumerate}

\section{Vì sao tập trung, và nó giải quyết gì}

Quy tắc twelve-factor là \emph{tách bạch nghiêm ngặt cấu hình khỏi
code}: một image được build một lần và thăng cấp nguyên vẹn qua các môi
trường; chỉ cấu hình bao quanh nó khác đi. Mọi cơ chế cấu hình đều là
một cách thỏa mãn quy tắc đó, và config-server là cơ chế đầu tiên hệ
sinh thái Spring sinh ra. Nó giải quyết ba vấn đề cụ thể cùng lúc:

\begin{itemize}
  \item \textbf{Trùng lặp.} Năm service dùng chung địa chỉ Kafka,
  endpoint tracing, mẫu log. Một tài liệu cho mỗi service ở một chỗ tốt
  hơn năm bản sao trôi dạt.
  \item \textbf{Nhận biết môi trường.} Cùng một service hỏi
  \code{/course-service/prod} hay \code{/course-service/default} và nhận
  đúng tài liệu, từ cùng một image.
  \item \textbf{Thay đổi không cần build lại.} Với backend git, một thay
  đổi cấu hình là một commit --- được review, có audit, hoàn tác được ---
  và có hiệu lực ở lần khởi động hay refresh kế tiếp, không build lại
  image nào.
\end{itemize}

Dự án nhận trọn lợi ích đầu, không nhận lợi ích thứ hai (chỉ có một
profile, \code{default}), và lợi ích thứ ba chỉ trên nguyên tắc: với
backend \code{native} các tài liệu nằm trong chính jar của
config-server, nên đổi một giá trị \emph{đúng là} phải build lại.

\begin{decision}[ba backend, ba hệ thống]
\textbf{Native classpath} (dự án này): cấu hình và code trong một repo,
một pipeline, một phiên bản. Đúng khi một đội sở hữu mọi thứ và mọi
thay đổi cấu hình đằng nào cũng đi cùng thay đổi code. \textbf{Backend
git}: cấu hình trong repository riêng với lịch sử riêng; đúng khi vận
hành và phát triển là những người khác nhau, khi cùng một service chạy
ở nhiều môi trường, hay khi kiểm toán viên hỏi ``ai đổi timeout và khi
nào''. \textbf{Nền tảng} (ConfigMap và Secret, không server): đúng khi
mọi đích triển khai đều là Kubernetes, vì nền tảng đã cung cấp cấu hình
có phiên bản, mount được, bảo vệ bằng RBAC và bớt một JVM phải chạy. Dự
án đang ở cái đầu và hướng tới cái thứ ba
(Mục~\ref{sec:config-enhance}); cái thứ hai là thứ phần lớn doanh
nghiệp bạn phỏng vấn thực sự đang chạy.
\end{decision}

\section{Đọc đường đi thật}

Theo dõi một thuộc tính, \code{spring.datasource.url} của
course-service, từ file tới pod đang chạy. Nó bắt đầu trong jar của
config-server:

\begin{lstlisting}[language=yaml]
# config-server/src/main/resources/configurations/course-service.yml
server:
  port: 8082
spring:
  datasource:
    url: jdbc:postgresql://localhost:5432/ts_course   (*@\co{1}@*)
    username: postgres
    password: root                                     (*@\co{2}@*)
app:
  s3:
    endpoint: ${APP_S3_ENDPOINT:http://localhost:9000}  (*@\co{3}@*)
    secret-key: ${APP_S3_SECRET_KEY:minioadmin}
\end{lstlisting}

File \code{application.yml} đóng gói của service nói lấy tài liệu đó ở
đâu:

\begin{lstlisting}[language=yaml]
# course-service/src/main/resources/application.yml
spring:
  application:
    name: course-service                                (*@\co{4}@*)
  config:
    import: optional:configserver:http://localhost:8888 (*@\co{5}@*)
\end{lstlisting}

Và container ghi đè cả địa chỉ lẫn giá trị:

\begin{lstlisting}[language=yaml]
# deploy/k8s/base/course-service.yaml (env excerpt)
- { name: SPRING_CONFIG_IMPORT,
    value: "configserver:http://config-server:8888" }    (*@\co{6}@*)
- { name: SPRING_DATASOURCE_URL,
    value: "jdbc:postgresql://postgres-course:5432/ts_course" }
envFrom:
  - secretRef: { name: postgres-course-app }            (*@\co{7}@*)
\end{lstlisting}

\begin{description}
  \item[\co{1}] Một giá trị cứng, không phải placeholder: tài liệu được
  phục vụ nói \code{localhost}. Nó chỉ đúng trên laptop.
  \item[\co{2}] Mật khẩu cơ sở dữ liệu dạng cleartext, trong một file
  git theo dõi, phục vụ qua HTTP. Ý định là chỉ cho dev; không gì trong
  file nói thế.
  \item[\co{3}] Thành ngữ \code{\${VAR:default}}: config-server phục vụ
  placeholder \emph{chưa phân giải}; mỗi client tự phân giải với môi
  trường của mình lúc khởi động. Đó là lý do cùng một tài liệu chạy
  được trong Compose và Kubernetes --- và là lý do một giá trị mặc định
  trong đó là giá trị mọi môi trường sẽ lặng lẽ dùng trừ khi được bảo
  khác.
  \item[\co{4}] Khóa tra cứu. Server ánh xạ \code{course-service} sang
  \code{course-service.yml}.
  \item[\co{5}] \code{optional:} --- nếu không gì trả lời ở 8888, cứ
  khởi động từ file đóng gói. Đúng cho một dev chạy một service trong
  IDE; sai cho một container.
  \item[\co{6}] Container thay cả dòng import: host mới, và
  \emph{không} có \code{optional:}. Config-server không với tới được giờ
  là chết người lúc khởi động. Compose che việc này bằng
  \code{depends\_on: condition: service\_healthy}; Kubernetes thì không,
  và pod crash-loop cho tới khi server Ready (Chương~15).
  \item[\co{7}] Mật khẩu ở \co{2} bị Secret ghi đè bằng
  \code{SPRING\_DATASOURCE\_PASSWORD}. Giá trị mặc định cleartext vì thế
  vô hại trong cụm --- \emph{chừng nào} Secret còn tồn tại. Xóa nó đi
  và service không thất bại; nó kết nối bằng \code{root}.
\end{description}

\subsection{Ưu tiên, chính xác}

Spring Boot dựng một \code{Environment} từ một danh sách có thứ tự các
nguồn thuộc tính; nguồn đầu tiên định nghĩa một khóa sẽ thắng. Với dự
án này thứ tự quan trọng, cao nhất trước:

\begin{enumerate}
  \item Tham số dòng lệnh và \code{SPRING\_APPLICATION\_JSON} (không
  dùng ở đây).
  \item Biến môi trường của OS --- \code{env} và \code{envFrom} của
  Deployment, với relaxed binding biến \code{SPRING\_DATASOURCE\_URL}
  thành \code{spring.datasource.url}.
  \item Tài liệu import từ config-server.
  \item File \code{application.yml} đóng gói đã thực hiện import.
\end{enumerate}

Hai hệ quả dễ trả lời sai trong phỏng vấn. Thứ nhất, tài liệu được
import đứng trên file đã import nó: config-server ghi đè được mọi thứ
đóng gói trong jar, nhưng không bao giờ ghi đè được thứ nền tảng đặt.
Thứ hai, placeholder \code{\${VAR:default}} \emph{không} phải quy tắc ưu
tiên; nó được phân giải sau khi ưu tiên đã chọn xong tài liệu. Biến môi
trường \code{APP\_S3\_SECRET\_KEY} thắng giá trị mặc định không phải vì
biến môi trường đứng trên config-server, mà vì chính văn bản của
config-server đã yêu cầu nó.

\begin{handson}[đọc thứ server phục vụ]
Config-server không có xác thực và Compose publish cổng của nó. Hỏi nó
tài liệu của course-service đúng như client làm:
\begin{lstlisting}[language=bash]
$ curl -s http://localhost:8888/course-service/default | \
    python -m json.tool | head -20
{
    "name": "course-service",
    "profiles": ["default"],
    "label": null,
    "version": null,
    "propertySources": [
        {
            "name": "classpath:/configurations/course-service.yml",
            "source": {
                "server.port": 8082,
                "spring.datasource.url": "jdbc:postgresql://localhost...",
                "spring.datasource.username": "postgres",
                "spring.datasource.password": "root",
\end{lstlisting}
Để ý \code{version: null} --- backend native không có commit id để báo.
Backend git điền vào đó mã SHA, và đó là cách bạn trả lời ``pod này
đang chạy cấu hình nào?'' trong production.
\end{handson}

\section{Ưu và nhược, thành thật}

\begin{center}
\begin{tabular}{@{}p{0.3\linewidth}p{0.62\linewidth}@{}}
\toprule
\textbf{Điểm mạnh} & \textbf{Trong dự án này} \\
\midrule
Một tài liệu cho mỗi service & Thật: năm file, không trùng lặp địa chỉ Kafka hay Zipkin. \\
Placeholder theo môi trường & Thật: cùng các file phục vụ Compose và Kubernetes không sửa gì. \\
Lấy lúc khởi động, không phụ thuộc lúc chạy & Thật: đã lên rồi, service không nói chuyện với server nữa. \\
\midrule
\textbf{Điểm yếu} & \\
\midrule
Không xác thực & Ai với tới 8888 đọc được mọi tài liệu, kể cả mặc định và mật khẩu. \\
Native = build lại để đổi & Chỉnh một timeout là một lượt build image config-server và rollout. \\
Mặc định dùng được cho secret & \code{JWT\_SECRET} và \code{root} có giá trị dự phòng dùng được trong git. \\
Ràng buộc thứ tự & Mọi pod phụ thuộc vào một pod không còn việc gì sau khi boot. \\
Chỉ một profile & \code{default} ở mọi nơi; môi trường chỉ khác nhau bằng biến môi trường. \\
\bottomrule
\end{tabular}
\end{center}

Dòng cuối đáng một câu. Profile của Spring được thiết kế đúng cho việc
tách dev/staging/prod, và config-server tự động phục vụ
\code{course-service-prod.yml} chồng lên \code{course-service.yml} khi
\code{SPRING\_PROFILES\_ACTIVE=prod}. Dự án không bao giờ dùng việc này;
mọi khác biệt môi trường là một biến môi trường. Đó là một đơn giản hóa
bảo vệ được --- một cơ chế thay vì hai --- nhưng nghĩa là các tài liệu
được phục vụ không bao giờ nói được ``giá trị này chỉ cho development''.

\section{Cái gì gãy}

\begin{pitfall}[khóa ký nằm trong repository]
Triệu chứng: không có, và đó chính là vấn đề. Mở \code{api-gateway.yml}
và \code{auth-service.yml} được phục vụ:
\begin{lstlisting}[language=yaml]
app:
  jwt:
    secret: ${JWT_SECRET:mySecretKeyThatIsAtLeast256Bits...}
\end{lstlisting}
Nguyên nhân: giá trị mặc định là một khóa HS256 \emph{hợp lệ}, commit
vào một repository công khai, và \code{create-secrets.sh} mặc định
\code{JWT\_SECRET} thành đúng chuỗi đó. Bất kỳ triển khai nào quên đặt
biến sẽ ký token bằng khóa cả internet đọc được --- và với HS256
(Chương~28) đọc được khóa nghĩa là đúc được token cho bất kỳ user nào.
Cách sửa có hai nửa: bỏ giá trị mặc định để biến thiếu thì thất bại lúc
khởi động (\code{\${JWT\_SECRET}} không có dấu hai chấm --- Spring từ
chối phân giải), và coi mọi \code{\${VAR:value}} trong tài liệu được
phục vụ như một câu hỏi: ``tôi có ổn không nếu production chạy với giá
trị này?'' Mật khẩu và khóa không bao giờ qua được bài kiểm tra đó;
cổng và hostname thường qua.
\end{pitfall}

\begin{pitfall}[refresh một service, hỏng tất cả]
Triệu chứng: sau khi ``xoay JWT secret'' bằng \code{POST
/actuator/refresh} trên auth-service, mọi request qua gateway trả về
401 cho tới khi gateway khởi động lại. Nguyên nhân: refresh là
\emph{theo từng instance}. Auth-service đã dựng lại các bean
\code{@RefreshScope} và giờ ký bằng khóa mới; gateway vẫn xác minh bằng
khóa cũ. Kể cả khi cả hai được refresh, không có cửa sổ chồng lấn ---
mọi token phát trước lần refresh trở nên vô hiệu ngay khoảnh khắc
gateway refresh. Một khóa đối xứng dùng chung giữa bên ký và bên xác
minh không thể hot-swap được; nó cần nghi thức dựa trên \code{kid} của
Chương~28. Quy tắc chung: refresh dành cho các giá trị mà bản cũ và bản
mới cùng tồn tại được (một feature flag, một mức log, một timeout),
không phải cho các giá trị hai bên phải thống nhất.
\end{pitfall}

\begin{pitfall}[khỏe mạnh nhưng sai]
Triệu chứng: config-server báo \code{UP}, service khởi động, giá trị
sai đang có hiệu lực. Nguyên nhân: health indicator của config-server
kiểm tra \emph{backend} có trả lời --- với native, là thư mục classpath
có tồn tại. Nó không nói gì về nội dung. Một lỗi chính tả ở tên khóa
YAML (\code{acess-key}) được phục vụ vui vẻ và bị client bỏ qua vui vẻ,
rồi client dùng giá trị mặc định của nó. Không bên nào log lỗi. Phòng
thủ nằm ở client: bind cấu hình vào một record
\code{@ConfigurationProperties} với \code{@Validated} và
\code{@NotBlank} trên các trường bắt buộc, để khóa thiếu hay sai chính
tả thất bại lúc khởi động với tên thuộc tính trong thông báo.
\end{pitfall}

\section{Chuyện gì gãy lúc 3 giờ sáng}

Pod config-server bị evict lúc 3 giờ sáng; không gì khác bị.

\begin{itemize}
  \item \textbf{Service đang chạy: không sao.} Mỗi cái đã lấy tài liệu
  một lần lúc boot và giữ trong bộ nhớ. Không đường request nào chạm
  tới config-server. Đây là tính chất biện minh cho một replica duy
  nhất.
  \item \textbf{Service khởi động lại trong lúc sự cố: crash loop.}
  Import của nó không optional; lấy thất bại; JVM thoát; kubelet khởi
  động lại với back-off lũy thừa (Chương~24). Khi config-server trở
  lại, lần khởi động kế tiếp thành công --- sau tối đa năm phút chờ
  back-off. Hệ thống hội tụ; nó không hội tụ nhanh.
  \item \textbf{Rolling update trong lúc sự cố: đứng.} Pod mới không bao
  giờ Ready, pod cũ tiếp tục phục vụ (\code{maxUnavailable} của
  Chương~15), và rollout chờ. Hành vi đúng, nhưng pipeline deploy sẽ
  báo timeout.
  \item \textbf{Khôi phục \code{optional:}: tệ hơn.} Pod khởi động lại
  boot từ giá trị mặc định đóng gói --- \code{localhost:5432},
  \code{localhost:9092} --- thành Ready (probe của nó là
  \code{/v3/api-docs}, không cần cơ sở dữ liệu), nhận lưu lượng, và
  thất bại mọi request. Sự cố trông như ``config-server chết'' giờ thành
  ``course-service trả 500'', cách nguyên nhân hai bước nhảy.
\end{itemize}

Retry với back-off là điểm giữa: giữ import bắt buộc, nhưng cho client
thử một lúc trước khi chết. Nó cần \code{spring-retry} và
\code{spring-boot-starter-aop} trên classpath, rồi:

\begin{lstlisting}[language=yaml]
spring:
  cloud:
    config:
      fail-fast: true          # unreachable server is fatal ...
      retry:
        max-attempts: 20       # ... after 20 tries
        initial-interval: 2000 # 2 s, then x1.5 each time
        multiplier: 1.5
        max-interval: 15000    # capped at 15 s: ~4 minutes total
\end{lstlisting}

So với để kubelet làm việc retry: khởi động lại JVM tốn hai mươi giây
nạp class mỗi lần và back-off tăng tới năm phút; retry trong tiến trình
không tốn gì và chặn trên ở mười lăm giây. Cả hai đều hội tụ; một cái
lãng phí một CPU mà chiếc máy hai nhân không có.

\section{Nâng cấp giải pháp}
\label{sec:config-enhance}

Hai nâng cấp, theo thứ tự một doanh nghiệp sẽ làm.

\subsection{Nâng cấp 1: backend git với giá trị mã hóa}

Giữ config-server, chuyển tài liệu của nó vào một repository, và ngừng
commit secret dạng cleartext. Phía server:

\begin{lstlisting}[language=yaml]
# config-server/src/main/resources/application.yml
server:
  port: 8888
spring:
  application:
    name: config-server
  profiles:
    active: git                                        (*@\co{1}@*)
  cloud:
    config:
      server:
        git:
          uri: https://git.example.com/ts/config-repo.git
          default-label: main                          (*@\co{2}@*)
          search-paths: '{application}'                (*@\co{3}@*)
          clone-on-start: true                         (*@\co{4}@*)
          username: ${CONFIG_GIT_USER}
          password: ${CONFIG_GIT_TOKEN}
encrypt:
  key: ${CONFIG_ENCRYPT_KEY}                           (*@\co{5}@*)
management:
  endpoints:
    web:
      exposure:
        include: health
\end{lstlisting}

\begin{description}
  \item[\co{1}] Profile chọn backend. \code{native} và \code{git} là hai
  cái có sẵn; Vault và JDBC cũng tồn tại.
  \item[\co{2}] Label là một branch hay tag. Ghim môi trường production
  vào \code{release-2026.09} và một thay đổi cấu hình là một lượt merge,
  không phải một lượt sửa.
  \item[\co{3}] Một thư mục cho mỗi service trong repository:
  \code{course-service/application.yml}. Cùng tên file như hôm nay, sâu
  hơn một thư mục.
  \item[\co{4}] Clone lúc khởi động thay vì lúc có request đầu tiên, để
  health check chứng minh credential dùng được trước khi client nào
  hỏi.
  \item[\co{5}] Một khóa đối xứng cho các endpoint \code{/encrypt} và
  \code{/decrypt}. Đó là secret duy nhất config-server phải giữ, và nó
  phải đến từ nền tảng (một Secret của Kubernetes), không bao giờ từ một
  file trong bất kỳ repository nào.
\end{description}

Có khóa rồi, mã hóa một giá trị là một request, và repository chỉ lưu
bản mã:

\begin{lstlisting}[language=bash]
$ curl -s -X POST http://localhost:8888/encrypt \
       -d 'a-real-256-bit-signing-key-generated-with-openssl'
AQBt3s6...b3f0e   # ciphertext; safe to commit
\end{lstlisting}

\begin{lstlisting}[language=yaml]
# config-repo/api-gateway/application.yml
app:
  jwt:
    secret: '{cipher}AQBt3s6...b3f0e'   # decrypted by the server on serve
\end{lstlisting}

Tiền tố \code{\{cipher\}} bảo server giải mã trước khi phục vụ.
Repository giờ không chứa secret dùng được nào, bẫy giá trị mặc định
biến mất (không có mặc định), và \code{git log -- api-gateway/} trả lời
ai đổi khóa và khi nào. Thứ nó \emph{không} sửa: tài liệu được phục vụ
vẫn đi qua mạng dạng cleartext tới một endpoint không xác thực. Thêm
\code{spring-boot-starter-security} vào config-server với credential
HTTP basic, và cấp cho mỗi client
\code{spring.cloud.config.username/password} từ Secret riêng của nó.

\subsection{Nâng cấp 2: cho server nghỉ hưu}

Trên Kubernetes nền tảng có thể tự phục vụ các tài liệu. Kustomize sinh
một ConfigMap cho mỗi service thẳng từ các file hiện có, nên lịch sử
git không đổi:

\begin{lstlisting}[language=yaml]
# deploy/k8s/base/kustomization.yaml (excerpt)
configMapGenerator:
  - name: course-service-config
    files:
      # key in the ConfigMap = application.yml, value = the file
      - application.yml=../../config-server/src/main/resources/configurations/course-service.yml
generatorOptions:
  disableNameSuffixHash: false   # hash suffix: a changed file = a
                                 # new ConfigMap name = a rollout
\end{lstlisting}

Deployment mount nó như một thư mục và trỏ Spring vào thư mục thay vì
server:

\begin{lstlisting}[language=yaml]
# deploy/k8s/base/course-service.yaml (excerpt)
spec:
  template:
    spec:
      containers:
        - name: course-service
          env:
            # replaces configserver:http://config-server:8888
            - { name: SPRING_CONFIG_IMPORT, value: "file:/config/" }
            - { name: SPRING_DATASOURCE_URL,
                value: "jdbc:postgresql://postgres-course:5432/ts_course" }
          envFrom:
            - secretRef: { name: postgres-course-app }
          volumeMounts:
            - { name: config, mountPath: /config, readOnly: true }
      volumes:
        - name: config
          configMap: { name: course-service-config }
\end{lstlisting}

Spring Boot coi dấu gạch chéo cuối là import thư mục và nạp
\code{application.yml} từ đó, với đúng thứ tự ưu tiên mà tài liệu từ
config-server từng có: trên file đóng gói, dưới biến môi trường. Không
gì trong code hay jar của service thay đổi; Compose vẫn chạy vì import
\code{optional:configserver:} đóng gói vẫn trỏ tới server (vẫn chạy,
trong Compose). Rồi xóa \code{config-server.yaml} khỏi base, và cụm bớt
một JVM, bớt một phụ thuộc khởi động, bớt một endpoint không xác thực.

Hậu tố hash trong tên là chi tiết khiến cách này tốt hơn server, chứ
không chỉ tương đương. Một ConfigMap tên
\code{course-service-config-7b9c4} đổi tên khi nội dung đổi; Deployment
tham chiếu tên mới là một pod template mới; Kubernetes roll các pod.
Thay đổi cấu hình, rollout và rollback là cùng một thao tác với thay đổi
code --- đúng thứ mục GitOps của Chương~27 muốn.

\begin{decision}[khi nào bỏ hẳn config-server]
Bỏ khi ba điều đúng: mọi runtime đều là Kubernetes, không service nào
cần giá trị đổi mà không khởi động lại pod, và secret đã sống trong nền
tảng (đúng thế --- Chương~16). Giữ khi bất kỳ điều nào sai: service
trên VM hay trong Compose ở production, một yêu cầu refresh lúc chạy
thật sự (hiếm; feature flag được phục vụ tốt hơn bởi hệ thống chuyên
dụng), hay một tổ chức mà dấu vết audit cấu hình phải độc lập với
repository triển khai. Người phỏng vấn hỏi ``bạn có dùng Spring Cloud
Config trên Kubernetes không?'' đang chờ nghe danh sách này, và câu
``nền tảng đã cung cấp sẵn''.
\end{decision}

\section{Ngoài dự án}

Config-server backend git với giá trị \code{\{cipher\}} là chuẩn từ
khoảng 2015 tới 2020 và vẫn là thứ bạn gặp ở phần lớn công ty Java. Hai
thứ đã thay thế từng phần của nó. HashiCorp Vault (và các bản tương
đương trên cloud, AWS Secrets Manager, Azure Key Vault) giữ secret với
lease và xoay vòng; Spring Cloud Vault hay External Secrets Operator
đưa chúng vào pod, và repository cấu hình chỉ giữ tham chiếu. Các dịch
vụ feature flag (LaunchDarkly, Unleash, OpenFeature) tiếp quản trường
hợp refresh lúc chạy với quy tắc nhắm mục tiêu và audit mà
\code{@RefreshScope} chưa bao giờ có. Thứ sống sót nguyên vẹn là kỷ
luật: cấu hình ngoài image, thứ tự ưu tiên bạn thuộc lòng, và không
secret dùng được trong bất kỳ file nào bạn commit.

\section{Câu hỏi từ phỏng vấn thật}

\subsection*{Q: Bạn phải xoay mật khẩu cơ sở dữ liệu cho năm mươi
service mà không có downtime. Trình bày từng bước.}

Xoay vòng không bao giờ là một lần ghi; nó là một sự chồng lấn. Thứ
nhất, tạo credential mới trong cơ sở dữ liệu trong khi cái cũ vẫn dùng
được --- Postgres cho phép hai role hay hai mật khẩu cùng tồn tại, và
các dịch vụ được quản lý (RDS, Cloud SQL) có ``dual password'' đúng cho
việc này. Thứ hai, đưa giá trị mới vào Secret của nền tảng (hoặc vào
repository cấu hình dưới dạng \code{\{cipher\}}) và kích hoạt các bên
tiêu thụ: trên Kubernetes là một lượt \code{rollout restart} cho mỗi
Deployment, việc mà mẹo hash tên ConfigMap làm hộ bạn; với Spring Cloud
Bus là một \code{POST /actuator/busrefresh} mà mọi instance nhận qua
Kafka. Thứ ba, theo dõi log kết nối của cơ sở dữ liệu cho tới khi không
session nào dùng role cũ, rồi thu hồi. Cạm bẫy của dự án ở Chương~16
--- ``đổi Secret rồi, không có gì xảy ra'' --- là thứ bước hai ngăn
chặn. Câu trả lời senior gọi tên sự chồng lấn một cách tường minh và
nói bước nào được phép thất bại: thu hồi, thứ bạn chỉ cần hoãn lại.

\subsection*{Q: Một thay đổi cấu hình phải tới mọi instance trong vòng
một phút. Khởi động lại thì quá chậm. Bạn làm gì?}

Phân biệt loại giá trị trước. Một mức log, một rate limit, một feature
flag: đánh dấu bean \code{@RefreshScope}, mở \code{/actuator/busrefresh},
thêm \code{spring-cloud-starter-bus-kafka} --- dự án đã có Kafka --- và
một POST phát sóng một \code{RefreshRemoteApplicationEvent}; mỗi
instance lấy lại tài liệu và dựng lại các bean refresh-scoped trong vài
giây. Một URL cơ sở dữ liệu, một địa chỉ bootstrap Kafka, một khóa ký:
\emph{đừng} refresh những thứ đó, vì connection pool, producer hay bên
xác minh giữ giá trị cũ trong các object không được dựng lại sạch sẽ;
khởi động lại thay vào đó. Nếu yêu cầu thực ra là ``flag cho thử
nghiệm sản phẩm'', hãy nói rằng một dịch vụ feature flag (Unleash,
LaunchDarkly) với SDK poll vài giây một lần là công cụ chuyên dụng, còn
\code{@RefreshScope} là cách chữa cháy mười năm tuổi.

\subsection*{Q: Config server là điểm lỗi đơn. Bạn xử lý thế nào?}

Bắt đầu bằng thu hẹp nó là điểm lỗi đơn \emph{cho cái gì}: khởi động,
không phải lưu lượng. Service đang chạy không bao giờ gọi lại nó, nên
sự cố config-server chỉ ảnh hưởng khởi động lại và scale-out. Rồi giảm
cả điều đó: hai replica sau Service (nó phi trạng thái với backend
native hay git --- mỗi replica một bản clone git là ổn),
\code{fail-fast} cộng retry để pod đang boot chờ tới vài phút thay vì
chết, và một readiness probe trên config-server để Service không bao
giờ định tuyến tới replica còn đang clone. Câu trả lời mạnh nhất loại
bỏ phụ thuộc: phục vụ tài liệu từ ConfigMap, và ``điểm lỗi đơn'' trở
thành API server của Kubernetes, thứ bạn đã phụ thuộc cho mọi thứ khác.

\subsection*{Q: Làm sao giữ secret ngoài git mà cấu hình vẫn trong git?}

Ba tầng, mỗi tầng chấp nhận được ở một quy mô khác nhau. Nhỏ nhất:
\code{/encrypt} của config-server với khóa đối xứng chỉ server giữ; git
lưu văn bản \code{\{cipher\}}, và repository bị lộ vô dụng nếu không có
khóa. Trung bình: Secret của Kubernetes tạo ngoài luồng
(\code{create-secrets.sh} của dự án) và được tham chiếu theo tên trong
manifest, nên git giữ hình dạng nhưng không bao giờ giữ giá trị. Lớn
nhất: một trình quản lý secret --- Vault, AWS Secrets Manager --- với
External Secrets Operator đồng bộ vào Secret của Kubernetes, để xoay
vòng và audit sống ở nơi secret sống. Cái bẫy cần gọi tên: một
\code{\${VAR:default}} với giá trị mặc định dùng được là một secret
trong git bất kể bạn chọn tầng nào.

\subsection*{Q: Hai pod của cùng một service đang chạy với cấu hình
khác nhau. Sao lại thế, và ngăn bằng cách nào?}

Nguyên nhân thường gặp: một rolling update kẹt nửa chừng (pod cũ với
tài liệu cũ, pod mới với tài liệu mới --- bình thường và tạm thời); một
config-server có bản clone git cũ trên một replica vì
\code{clone-on-start} tắt và replica đó không bao giờ làm mới; hay,
tinh vi nhất, một \code{/actuator/refresh} gửi tới một instance qua load
balancer tình cờ chọn nó. Phòng ngừa là biến cấu hình thành một phần
của pod template: mount từ ConfigMap có tên mang hash nội dung, để
``cấu hình nào'' được trả lời bằng ``ReplicaSet nào''. Phát hiện là
phơi phiên bản cấu hình --- mã SHA git mà config-server trả về dưới
dạng \code{version}, hoặc tên ConfigMap --- trong \code{/actuator/info},
và so sánh giữa các pod trước khi tin bất kỳ báo cáo ``nó chạy trên
một trong số chúng'' nào.

\begin{quiz}
\begin{enumerate}
  \item \code{course-service.yml} được phục vụ nói \code{password: root}
  và Secret của Deployment nói khác. Nêu cái nào thắng và vì sao, rồi mô
  tả chuyện gì xảy ra nếu Secret bị xóa và pod khởi động lại.
  \item Giải thích vì sao \code{\${APP\_S3\_SECRET\_KEY:minioadmin}}
  trong tài liệu được phục vụ không phải quy tắc ưu tiên, và
  config-server thực sự gửi gì cho client với khóa đó.
  \item Một đồng nghiệp chạy \code{POST /actuator/refresh} trên cả
  auth-service và api-gateway sau khi đổi \code{JWT\_SECRET}. Liệt kê
  hai cửa sổ lỗi riêng biệt và nêu cơ chế từ Chương~28 mà xoay khóa
  nóng thực sự cần.
  \item Config-server chết; course-service được scale từ một lên hai
  replica trong lúc sự cố. Mô tả pod mới làm gì với \code{configserver:}
  (không optional), với \code{optional:configserver:}, và với
  \code{fail-fast} cộng retry.
  \item Sau khi chuyển sang ConfigMap, cái gì khiến pod khởi động lại
  khi một giá trị đổi, và vì sao điều đó khiến việc chuyển đổi \emph{có
  năng lực hơn} server nó thay thế chứ không chỉ tương đương?
  \item Phỏng vấn: ``Bạn có triển khai Spring Cloud Config trên
  Kubernetes không?'' Nêu ba điều kiện để bạn bỏ nó và một điều kiện để
  bạn giữ.
\end{enumerate}
\end{quiz}
